%Drehstrom
\begin{frame}
    \fta{Three-phase current}
    
    \s{
        To ensure the most efficient energy transport possible, three-phase current is used in many grid systems. 
        Three-phase current is based on alternating current, which is not taken as a single phase, but as a combination of several phases.
        Alternating current is generated by generators that have not just a single winding in the rotor, but several.       
        If these windings are installed at a symmetrical distance from each other, phases are generated depending on the number of windings.
        This is then referred to as a multi-phase system. If there are exactly three phases, the system is called three-phase current.
        The basics of alternating current and complex calculations are particularly relevant for understanding three-phase current.
    }
    
    \begin{Learning objectives}{Drehstrom}
        The students
        \begin{itemize}
            \item understand the basic characteristics of three-phase current and the associated operators.
            \item can analyse three-phase systems with the same load (symmetrical components).
            \item are familiar with the characteristics of three-phase systems with different phase loads (unbalanced
            components).
        \end{itemize}
    \end{Learning objectives}
    
\end{frame}

\begin{frame}
    \ftb{Symmetrical components}
    
    
    \s{
        In the field of energy generation, transmission and distribution, three-phase current is mainly used.
        The term ‘three-phase current’ comes from the technical fact that the systems are connected to the spatially 
        rotating magnetic field of the generator: the rotating field.    
        All phases in a multi-phase system are generated at the same frequency, which in Europe is known to be 50 Hz.
        The generated phases are connected at the point of generation, i.e. the generator, and behind the consumers.
        In a three-phase system, star connection or delta connection are available for this purpose.   
        In the symmetrical case, the generator generates as many phases as there are windings in the rotor.
        The windings are always installed at the same angle. This angle is assigned the Greek letter $\alpha$:
    }
    
    \b{
        \begin{itemize}
            \item Phases are generated depending on the number of windings
            \item In the symmetrical case, the windings are installed in the generator with the same angle $\alpha$
        \end{itemize}
    }
    
    
    \begin{eqa}
        \alpha = \frac{2\pi}{m}
    \end{eqa}
    
    
    \b{
        \begin{itemize}
            \item with $m$ = number of phases
            \item For a three-phase system: $m=3$
        \end{itemize}
        \begin{eqa}
            \alpha=\frac{2\pi}{3}=120 \degree
        \end{eqa}
    }
    
    \s{
        In the formula, m stands for the number of phases generated in the system.
        In a three-phase system, there are three phases.
        This results in an $\alpha$ of 120°.
    }
    
\end{frame}

\begin{frame}
    \renewcommand{\figurename}{Figure}
    \ftx{Symmetrical components}
    
    \s{
        \fu{
            %vergrößern auf Folienbreite
            \resizebox{0.8\textwidth}{!}{
                \input{Tikz/src/Drehstrom/Drei_Phasen_Drehstrom}
            }
        }{{\bf Three phases of three-phase current.} The three phases of the three-phase system are all phase-shifted by 120° relative to each other.
        This results in an angle $\alpha$ of 120°.}
        
        
        The next step is to derive a rotation operator from the angle $\alpha$.
        This rotation operator is used to convert the voltages and currents to the reference voltage/current.
    }
    
    \b{
        The rotation operator a as a function of $\underline{\alpha}$:
    }
    
    
    \begin{eqa}
        \underline{a}=e^{j\alpha}=e^{j\frac{2\pi}{m}} \label{Drehoperator}
    \end{eqa}
    
    For $m=3$:
    
    \begin{eqa}
        \underline{a}&=e^{j\frac{2\pi}{3}}=-\frac{1}{2}+j\frac{\sqrt{3}}{2}                         \notag \\      
        \underline{a}^2&=e^{j\frac{4\pi}{3}}=-\frac{1}{2}-j\frac{\sqrt{3}}{2}=\underline{a}^{*}     \notag \\  
        \underline{a}^3&=e^{j2\pi}=1                                                                \notag \\
        \underline{a}^4&=e^{j\frac{8\pi}{3}}=-\frac{1}{2}+j\frac{\sqrt{3}}{2}=\underline{a}         \notag \\
        \underline{a}^5&=e^{j\frac{10\pi}{3}}=-\frac{1}{2}-j\frac{\sqrt{3}}{2}=\underline{a}^2      \notag \\
        \vdots  \notag
    \end{eqa}
\end{frame}

\begin{frame}
    \ftx{Rotation operator}
    
    \begin{Key point}{Rotation operator}
        The angle $\alpha$ indicates the phase shift between the phases. With the help of the rotation operator $\underline{a}$,
        voltages and currents can be converted. 
    \end{Key point}
\end{frame}

\begin{frame}
    \ftx{Symmetrical components}
    \s{
        The complex voltages of the respective outer conductors relative to a neutral conductor can be formulated using the following equation:    }

    \b{
        The following applies to the voltage of the respective outer conductors relative to a neutral conductor:
    }
    
    \begin{eqa}
        \underline{U}_{\mathrm{Lm}}=U_\mathrm{L1}\cdot e^{-j(m-1)\alpha}
    \end{eqa}
    
    \s{
        The voltage $U_1$ is a real reference value for the multiphase system.

        Now the equations for the rotation operator are combined with those for the voltage, resulting in the following relationship:
    }
    
    \b{
        Voltage $U_m$ as a function of the rotation operator $a$:
    }
    
    \begin{eqa}
        \underline{U}_\mathrm{L1}&=\underline{a}^m\cdot U_\mathrm{L1}=1\cdot U_\mathrm{L1}         \notag \\       
        \underline{U}_\mathrm{L2}&=\underline{a}^{m-1}\cdot U_\mathrm{L1}                \notag \\
        \vdots                                                       \notag \\
        \underline{U}_{\mathrm{Lm-1}}&=\underline{a}^2\cdot U_\mathrm{L1}       \notag \\
        \underline{U}_{\mathrm{Lm}}&=\underline{a}\cdot U_\mathrm{L1}           \notag
    \end{eqa}
    
\end{frame}

\begin{frame}
    \ftx{Symmetrical components}
    
    With three phases, i.e. $m=3$, the following relationship between the three voltages and the rotation operator results:
    
    \begin{eqa}
        \underline{U}_\mathrm{L1}+\underline{U}_\mathrm{L2}+\underline{U}_\mathrm{L3}=U_1(1+\underline{a}^2+\underline{a})
    \end{eqa}
    
    \s{
        \textit{Note:}
        
        In the literature, the indices for the stresses are not always chosen uniformly. The most common alternative is the letters $u, v, w$ for phases $1, 2, 3$.
        Other indices are also used, but these will not be discussed further here.
    }
    
\end{frame}

\begin{frame}
    \renewcommand{\figurename}{Figure}
    \ftx{Symmetrical components}
    
    \s{
        Fig. \ref{ZeigerDrehoperator} illustrates once again how the rotation operators are composed in the complex numbers.
    }
    
    \fu{
        \resizebox{0.8\textwidth}{!}{%
            \input{Tikz/src/Drehstrom/Zeigerdiagramm_Drehstromoperator}
        }
    }{{\bf Pointer diagram of the rotation operator in the three-phase system.} If the rotation operators are added together, the result is
        zero. \label{ZeigerDrehoperator}}
    
    \b{
        Pointer diagram for the rotation operator $\rightarrow$ The sum is zero.
        \begin{eqa}
            1+\underline{a}+\underline{a}^2=0                       \notag \\
            \underline{U}_\mathrm{L1}+\underline{U}_\mathrm{L2}+\underline{U}_\mathrm{L3}=0       \notag
        \end{eqa}
        
    }
    
    \s{
        If you follow the path via the three rotation operators, you return to the origin and obtain zero.

        \begin{eqa}
            1+\underline{a}+\underline{a}^2=0                       \notag \\
            \underline{U}_\mathrm{L1}+\underline{U}_\mathrm{L2}+\underline{U}_\mathrm{L3}=0       \notag
        \end{eqa}
    }
    
\end{frame}


\begin{frame}
    \ftc{Line quantities in a star (wye) connection}
    
    \s{
        The following chapters will discuss the types of connections that can occur in a three-phase system.
        First, we will look at what is known as a star connection. 
        A star connection means that the phases are connected in such a way that a so-called star point is formed.
        A system can be connected on both the generator and consumer sides.
        Fig. \ref{ESBStern} shows the typical components of a star connection in a three-phase equivalent circuit diagram.
        A voltage source is entered for each phase generated by the generator.
        If the generator has a star connection, the individual voltages are combined in a node in the ESB.
        This node is designated by the letter $N_Q$.         
        A strand emanates from each voltage source, which should reflect the respective phase.
        Due to physical conditions, the components connected downstream of the generator, i.e. primarily the line and the consumer, exhibit both active and reactive resistance.
        These are combined in the ESB to form one impedance per strand.       
        The impedances of the strands are combined in the star connection, as are the voltage sources, to form the node $N_V$.
        In a system in which both the generator and consumer sides are connected in a star configuration, there may be a neutral conductor or return conductor in addition to the phase strands, which connects the two nodes.
    }
    
    \b{
        \begin{itemize}
            \item The generator and consumer sides can be connected either in a star or a delta configuration
            \item In a star connection, the phases are connected to a star point
            \item There are then nodes on the generator and consumer sides that can be connected via a neutral conductor
            \item A voltage source is entered for each phase generated by the generator
            \item The consumers connected to each phase are combined into one impedance.
        \end{itemize}
    }
    
\end{frame}


\begin{frame}
    \renewcommand{\figurename}{Figure}
    \ftx{Interlinked systems in a star configuration}
    
    \fu{
        \resizebox{0.8\textwidth}{!}{%
            \input{Tikz/src/Drehstrom/ESB_Drehstromsystem}
        }
    }{{\bf ESB for a three-phase system.} Generators and consumers in the star and a neutral conductor between the nodes. \label{ESBStern}}
    
    \b{
        ESB for a three-phase system with generators and consumers in a star configuration and return conductor
    }
\end{frame}

\begin{frame}
    \ftx{Interlinked systems in a star configuration}
    
    
    \textbf{Voltages}
    
    In principle, there are three voltages in a three-phase system that should be specifically mentioned:

    \begin{itemize}
        \item Phase-to-neutral voltage (star voltage) $U_\Stern$: Voltage of the conductor relative to the neutral conductor ($U_{\mathrm{L1}}$, $U_{\mathrm{L2}}$, $U_{\mathrm{L3}}$)
        \item External conductor voltage (chained voltage) U: Voltage between two conductors ($U_{\mathrm{L1L2}}$, $U_{\mathrm{L2L3}}$, $U_{\mathrm{L3L1}}$) -- this voltage is also referred to as delta voltage
        \item Phase voltage $U_{\mathrm{str}}$: Voltage across the impedance of the conductor relative to the neutral conductor. Depending on the connection, the phase voltage is either equal to the star voltage or the delta voltage.
    \end{itemize}
    
\end{frame}


\begin{frame}
    \ftx{Interlinked systems in a star configuration}
    
    \b{
    \textbf{Currents}
    
    The current in the neutral conductor is determined by Kirchhoff's 1st law:
    \begin{eqa}
        \underline{I}_{\mathrm{N}}=\underline{I}_{\mathrm{L1}}+\underline{I}_{\mathrm{L2}}+\underline{I}_{\mathrm{L3}}
    \end{eqa}
    
    Converted according to Ohm's law:
    \begin{eqa}
        \underline{I}_{\mathrm{N}}=\frac{\underline{U}_{\mathrm{L1}}}{\underline{Z}_1}+\frac{\underline{U}_{\mathrm{L2}}}{\underline{Z}_2}+\frac{\underline{U}_{\mathrm{L3}}}{\underline{Z}_3}=\underline{U}_{\mathrm{1m}}\underline{Y}_{1}+\underline{U}_{\mathrm{2m}}\underline{Y}_{2}+\underline{U}_{\mathrm{3m}}\underline{Y}_{3}
    \end{eqa}
    
    In the symmetrical case, all impedances are equal ($\underline{Z}_{\mathrm{L1}}=\underline{Z}_{\mathrm{L2}}=\underline{Z}_{\mathrm{L3}}=\underline{Z}$), from which it follows that:
    \begin{eqa}
        \underline{I}_{\mathrm{N}}=\frac{1}{\underline{Z}}(\underline{U}_{\mathrm{L1}}+\underline{U}_{\mathrm{L2}}+\underline{U}_{\mathrm{L3}})=0
    \end{eqa}
    
    }
    
    \s{
        \textbf{Currents}
        
        For currents in a star system, it is helpful to apply Kirchhoff's first law.
        The sum of the currents in the strands must be equal to the current in the neutral conductor:
        
        \begin{eqa}
            \underline{I}_{\mathrm{N}}=\underline{I}_{\mathrm{L1}}+\underline{I}_{\mathrm{L2}}+\underline{I}_{\mathrm{L3}}
        \end{eqa}
        
        In relation to voltage and impedance, Ohm's law states that:
        
        \begin{eqa}
            \underline{I}_{\mathrm{N}}=\frac{\underline{U}_{\mathrm{L1}}}{\underline{Z}_1}+\frac{\underline{U}_{\mathrm{L2}}}{\underline{Z}_2}+\frac{\underline{U}_{\mathrm{L3}}}{\underline{Z}_3}=\underline{U}_{\mathrm{1m}}\underline{Y}_{1}+\underline{U}_{\mathrm{2m}}\underline{Y}_{2}+\underline{U}_{\mathrm{3m}}\underline{Y}_{3}
        \end{eqa}
        
        When referring to a symmetrical system, this applies not only to the generator but also to the impedances.
        Accordingly, the impedances of the three strands are all equal:
        
        \begin{eqa}
            \underline{Z}_{\mathrm{L1}}=\underline{Z}_{\mathrm{L2}}=\underline{Z}_{\mathrm{L3}}=\underline{Z} \label{SymImpedanzen}
        \end{eqa}
        
        If this relationship is taken into account for a symmetrical system for equation \ref{SymImpedanzen}, the impedance can be extracted.
        It is already known that the sum of the voltages in a symmetrical system is zero.
        It can therefore be proven that the sum of the currents must also be zero.
        
        \begin{eqa}
            \underline{I}_{\mathrm{N}}=\frac{1}{\underline{Z}}(\underline{U}_{\mathrm{L1}}+\underline{U}_{\mathrm{L2}}+\underline{U}_{\mathrm{L3}})=0
        \end{eqa}
    }
    
\end{frame}

\begin{frame}
    \renewcommand{\figurename}{Figure}
    \ftx{Interlinked systems in a star configuration}
    
    
    
    \fu{
        \resizebox{0.9\textwidth}{!}{%
            \input{Tikz/src/Drehstrom/Dreisphasiges_ESB}
        }
    }{{\bf Three-phase equivalent circuit diagram star-star.} ESB with symmetrical loads in star configuration \label{ESBSymStern}}
    
    \b{
        Three-phase ESB with symmetrical consumers in star configuration
    }
    
\end{frame}


\begin{frame}
    \renewcommand{\figurename}{Figure}
    \ftx{Interlinked systems in a star configuration}
    
    \s{
        To avoid having to draw the complete three-phase ESB as shown in Fig. \ref{ESBSymStern}, it is possible to draw a simplified single-phase ESB.   
        When drawing the single-phase ESB (and, of course, any other ESB), it is important to use the correct designations for the currents, voltages and impedances via the indices.
        The single-phase ESB for the star connection is shown in Fig. \ref{EinphasigESBStern}.
    }
    
    \fu{
        \input{Tikz/src/Drehstrom/Einphasiges_ESB}
    }{{\bf Single-phase equivalent circuit diagram star-star.} ESB with symmetrical loads in star configuration \label{EinphasigESBStern}}
    
    \b{
        \begin{itemize}
            \item Alternatively, the ESB can also be set up as single-phase
            \item It is important to correctly label the currents, voltages, and impedances
        \end{itemize}
    }
    
\end{frame}

\begin{frame}
    \ftx{Symmetrical components}
    
    \begin{Key point}{Symmetrical components}
        {\bf Symmetrical components} describe an {\bf equal load on the phases}, 
        for example in a three-phase system. Here, all three phases are loaded by identical consumers.  
    \end{Key point}
\end{frame}


\begin{frame}
    \ftc{Linked systems in a triangle}
    
    \s{
        
        As an alternative to star connection, the sources and consumers can be connected in a delta configuration. 
        However, in most cases, the generator side is connected in a star configuration.
        The reason for this is that even the smallest asymmetries cause compensating currents to flow, which would affect the operation of the generator.      
        This should be avoided if possible, which is why a star connection is a better operating mode.
        For this reason, it is assumed that the generator side is fundamentally connected in a star configuration.
        So when we talk about a delta connection, unless explicitly stated otherwise, we are only referring to the consumer side.
        
        It has already been defined that the outer conductor voltage is also called the delta voltage.
        According to Mesh Law, the delta voltage of two conductors is calculated from the difference between the respective star voltages:
        
        \begin{eqa}
            \underline{U}_{\mathrm{L1L2}}=\underline{U}_{\mathrm{L1}}-\underline{U}_{\mathrm{L2}} \label{AußenleiterKomplex}
        \end{eqa}
        
        If the value of the respective outer conductor voltage is to be calculated, complex calculations would be required according to equation \ref{AußenleiterKomplex}.
        However, a trigonometric derivation can be used to describe the relationship between star voltage and delta voltage.    
        }
    \b{
        \begin{itemize}
            \item In a delta connection, the phases are connected directly to each other
            \item However, in most cases, the generator side remains connected in a star configuration in order to better dissipate compensation currents
            \item The outer conductor voltage (also known as delta voltage) is calculated from the difference between the respective star voltages:
        \end{itemize}
        
        \begin{eqa}
            \underline{U}_{\mathrm{L1L2}}=\underline{U}_{\mathrm{L1}}-\underline{U}_{\mathrm{L2}} 
        \end{eqa}
    }
    
\end{frame}

\begin{frame}
    \ftx{Interlinked systems in a triangle}
    
    \b{
        \begin{itemize}
            \item The relationship between triangular and star tension can be described using a trigonometric derivation:        \end{itemize}
    }
    
    \begin{eq}
        \sin\left(\frac{\alpha}{2}\right)=\sin\left(\frac{\frac{2\pi}{m}}{2}\right)=\sin\left(\frac{\pi}{m}\right)=\frac{G}{H} \label{GleichungW3}
    \end{eq}
    
    \begin{eqa}
        \sin\left(\frac{\pi}{3}\right)=\frac{\frac{\underline{U}_{\mathrm{L1L2}}}{2}}{\underline{U}_{\mathrm{L1}}} \notag 
    \end{eqa}
    \begin{eq}
        2\cdot\frac{\sqrt{3}}{2}=\frac{\underline{U}_{\mathrm{L1L2}}}{\underline{U}_{\mathrm{L1}}} \label{GleichungW32}
    \end{eq}
    \begin{eqa}
        \underline{U}_{\mathrm{L1}}=\frac{\underline{U}_{\mathrm{L1L2}}}{\sqrt{3}}     \notag
    \end{eqa}
    
\end{frame}


\begin{frame}
    \ftx{Star-delta conversion factor}
    
    \begin{Key point}{Stern-Dreieck-Umrechnungsfaktor}
        The star voltage and the delta voltage can be converted into each other using the conversion factor $\sqrt{3}$: 
        \begin{eq}
            U_{\Stern}=\frac{U_{\Dreieck}}{\sqrt{3}}    \nonumber
        \end{eq}
    \end{Key point}
\end{frame}


\begin{frame}
    \renewcommand{\figurename}{Figure}
    \ftx{Interlinked systems in a triangle}
    
    \s{
        In Fig. \ref{HerrleitungW3}, the phasor diagram of the complex rotating voltages is used to illustrate the relationship 
        between star and delta voltages. For this purpose, trigonometric principles are used 
        to represent the length of the outer conductor voltage (i.e. the delta voltage) in a different way. Equation \ref{GleichungW3}
        first describes in general terms how sine can be used in a multi-phase system. This equation 
        applies in general, regardless of how many phases the system has. In equation \ref{GleichungW32}, $m$ is set to 3.
        This means that the angle is always half the actual angle between the phases.
        This means that the opposite side of the sine is half the delta voltage. By rearranging and converting the value 
        $sin(\frac{\pi}{3})$ to $\frac{\sqrt{3}}{2}$, we finally obtain the factor needed to convert star and delta voltage.
    }
    
    \b{
        \begin{itemize}
            \item The derivation can be illustrated even better on the basis of the pointer diagram:
        \end{itemize}
    }
    
    \s{
        In practice, star voltage is always marked with a star symbol as an index, while delta voltage 
        usually does not require an index.
        
        \begin{eqa}
            U_{\Stern}=\frac{U_{\Dreieck}}{\sqrt{3}}=\frac{U}{\sqrt{3}}
        \end{eqa}
    }
    
    \fu{
    \resizebox{0.35\textwidth}{!}{%
        \input{Tikz/src/Drehstrom/Zeigerdiagramm_Herleitung_Stern_Dreieck}
    }
    }{{\bf Pointer diagram illustrating the derivation of the factor between star and delta connection.}
    The series voltage between two outer conductors $\underline{U}_{\mathrm{L1L2}}$ can be converted to the outer conductor voltage $\underline{U}_{\mathrm{L1}}$ using the 
    conversion factor $\sqrt{3}$.
    \label{HerrleitungW3}}
    
    \b{
        \begin{itemize}
            \item In practice, the triangle symbol is usually omitted from the index:
                  %$U_{\Stern}=\frac{U_{\Dreieck}}{\sqrt{3}}=\frac{U}{\sqrt{3}}$
        \end{itemize}
        
        \begin{eqa}
            U_{\Stern}=\frac{U_{\Dreieck}}{\sqrt{3}}=\frac{U}{\sqrt{3}}
        \end{eqa}
    }
\end{frame}

\begin{frame}
    \ftx{Interlinked systems in a triangle}
    
    \s{
        As can be seen in Figure \ref{ESBDreieck}, the impedances of the consumers are now connected in such a way that 
        there is no longer a neutral point. In this circuit, the phase voltages now correspond to the 
        outer conductor or delta voltage. 
    }
    
    \b{
        Three-phase ESB with generators in star configuration and consumers in delta configuration
    }
    
    \fu{\resizebox{\textwidth}{!}{%
            \input{Tikz/src/Drehstrom/Stern-Dreieck_dreiphasiges_ESB}
        }
    }{{\bf Three-phase equivalent circuit diagram star-delta.} ESB with producers in the star and consumers in the triangle. \label{ESBDreieck}}
\end{frame}


\begin{frame}
    \ftx{Interlinked systems in a triangle}
    
    
    \s{
        The currents between the strands can be determined from the node rule as follows:
        
        \begin{eqa}
            \underline{I}_{\mathrm{L1}}=\underline{I}_{\mathrm{L1L2}}-\underline{I}_{\mathrm{L3L1}}  \\
            \underline{I}_{\mathrm{L2}}=\underline{I}_{\mathrm{L2L3}}-\underline{I}_{\mathrm{L1L2}}  \\
            \underline{I}_{\mathrm{L3}}=\underline{I}_{\mathrm{L3L1}}-\underline{I}_{\mathrm{L2L3}} 
        \end{eqa}
        
        As with the star connection, in the symmetrical case it can be said that all impedances have the same value.
        This results in the following for the currents in the phase:
        
        \begin{eqa}
            \underline{I}_{\mathrm{L1}}&=\underline{I}_{\mathrm{L1L2}}-\underline{I}_{\mathrm{L3L1}}=\frac{\underline{U}_{\mathrm{L1L2}}-\underline{U}_{\mathrm{L3L1}}}{\underline{Z}}=\frac{U_{\Stern}}{\underline{Z}}\cdot (1-\underline{a}^2-\underline{a}+1)=3\cdot \frac{U_{\Stern}}{\underline{Z}} \\ 
            \underline{I}_{\mathrm{L2}}&=\underline{I}_{\mathrm{L2L3}}-\underline{I}_{\mathrm{L1L2}}=\frac{\underline{U}_{\mathrm{L2L3}}-\underline{U}_{\mathrm{L1L2}}}{\underline{Z}}=\frac{U_{\Stern}}{\underline{Z}}\cdot (\underline{a}^2-\underline{a}-1+\underline{a}^2)=3\cdot \underline{a}^2 \cdot \frac{U_{\Stern}}{\underline{Z}} \\
            \underline{I}_{\mathrm{L3}}&=\underline{I}_{\mathrm{L3L1}}-\underline{I}_{\mathrm{L2L3}}=\frac{\underline{U}_{\mathrm{L3L1}}-\underline{U}_{\mathrm{L2L3}}}{\underline{Z}}=\frac{U_{\Stern}}{\underline{Z}}\cdot (\underline{a}-1-\underline{a}^2+\underline{a})=3\cdot \underline{a} \cdot \frac{U_{\Stern}}{\underline{Z}}
        \end{eqa}
        
        
        These transformations show that the currents differ only in terms of the rotation operator $a$. 
        Equivalent to the conversion from star to delta voltage, a ratio between phase and delta current can also be established.
        Using the same derivation as shown in Fig. \ref{HerrleitungW3}, but replacing the voltage with the corresponding current values, the following ratio is obtained:        
        \begin{eqa}
            I_{\mathrm{str}}=\sqrt{3}\cdot I_{\Dreieck}
        \end{eqa}
        
        With the available information, a single-phase ESB can now also be set up for the delta connection (see Fig. \ref{EinphasigESBDreieck}).
    }
    
    \b{
        \begin{itemize}
            \item In a delta connection, the phase voltage corresponds to the outer conductor (delta) voltage
            \item In the symmetrical case, the impedances are also equal in this circuit
            \item The following equations can be established for the current:
        \end{itemize}
        \begin{eqa}
            \underline{I}_{\mathrm{L1}}=\underline{I}_{\mathrm{L1L2}}-\underline{I}_{\mathrm{L3L1}}  \\
            \underline{I}_{\mathrm{L2}}=\underline{I}_{\mathrm{L2L3}}-\underline{I}_{\mathrm{L1L2}}  \\
            \underline{I}_{\mathrm{L3}}=\underline{I}_{\mathrm{L3L1}}-\underline{I}_{\mathrm{L2L3}} 
        \end{eqa}
        \begin{eqa}
            \underline{I}_{\mathrm{L1}}&=\underline{I}_{\mathrm{L1L2}}-\underline{I}_{\mathrm{L3L1}}=\frac{\underline{U}_{\mathrm{L1L2}}-\underline{U}_{\mathrm{L3L1}}}{\underline{Z}}=\frac{U_{\Stern}}{\underline{Z}}\cdot (1-\underline{a}^2-\underline{a}+1)=3\cdot \frac{U_{\Stern}}{\underline{Z}} \\ 
            \underline{I}_{\mathrm{L2}}&=\underline{I}_{\mathrm{L2L3}}-\underline{I}_{\mathrm{L1L2}}=\frac{\underline{U}_{\mathrm{L2L3}}-\underline{U}_{\mathrm{L1L2}}}{\underline{Z}}=\frac{U_{\Stern}}{\underline{Z}}\cdot (\underline{a}^2-\underline{a}-1+\underline{a}^2)=3\cdot \underline{a}^2 \cdot \frac{U_{\Stern}}{\underline{Z}} \\
            \underline{I}_{\mathrm{L3}}&=\underline{I}_{\mathrm{L3L1}}-\underline{I}_{\mathrm{L2L3}}=\frac{\underline{U}_{\mathrm{L3L1}}-\underline{U}_{\mathrm{L2L3}}}{\underline{Z}}=\frac{U_{\Stern}}{\underline{Z}}\cdot (\underline{a}-1-\underline{a}^2+\underline{a})=3\cdot \underline{a} \cdot \frac{U_{\Stern}}{\underline{Z}}
        \end{eqa}
        
    }
\end{frame}


\begin{frame}
    \renewcommand{\figurename}{Figure}
    \ftx{Interlinked systems in a triangle}
    
    \b{
        \begin{itemize}
            \item The ratio of phase current to delta current can be calculated in the same way as the ratio of star voltage to delta voltage:
        \end{itemize}
        \begin{eqa}
            I_{\mathrm{str}}=\sqrt{3}\cdot I_{\Dreieck}
        \end{eqa}
        
        \begin{itemize}
            \item A single-phase ESB can also be set up for the delta connection:
        \end{itemize}
    }
    \fu{
        \input{Tikz/src/Drehstrom/Stern-Dreieck_einphasiges_ESB}
    }{{\bf Single-phase equivalent circuit diagram star-delta.} ESB of a delta connection with symmetrical consumers \label{EinphasigESBDreieck}}
    
\end{frame}




\begin{frame}
    \ftc{Power in a symmetrical three-phase network}
    \s{
        As in all electrical systems, power is calculated using the equation $P=U\cdot I$.
        In the case of a three-phase system, the power must be calculated per phase.
        In a symmetrical system, the currents and voltages per phase are equal, and the result for one phase must therefore be multiplied by 3.
        
        \begin{eqa}
            S&=3\cdot I_{\mathrm{str}}\cdot U_{\Stern}=3\cdot I_{\mathrm{str}}\cdot \frac{U_{\Dreieck}}{\sqrt{3}} \notag \\
            &=3\cdot I_{\Dreieck}\cdot U_{\Dreieck}=3\cdot \frac{I_{\mathrm{str}}}{\sqrt{3}}\cdot U_{\Dreieck} \notag \\
            &=\sqrt{3}\cdot U_{\Dreieck}\cdot I_{\mathrm{str}} \label{LeistungDrehstrom}
        \end{eqa}
         
        The formula \ref{LeistungDrehstrom} shows that it makes no difference to the power calculation which circuit is used.
        The only thing to bear in mind is which voltages and currents are used in the calculation (delta or star).
        
        \textit{Note:}
        
        The notation of the indices for delta or star values is not clear in the literature.
        In most cases, no index is given for the delta voltage (external conductor voltage) and the phase current $U_{\delta}=U$ and $I_{ph}=I$, but only for the star voltage and the delta current.        
    
        If you directly compare the performance of a delta connection and a star connection, you get different results.

        The following value is obtained for the star connection:
        \begin{eqa}
            S=\sqrt{3}\cdot U\cdot I=\sqrt{3}\cdot U \cdot \frac{U_{\Stern}}{Z}=\frac{U^2}{Z}
        \end{eqa}
        For the delta connection, we again obtain:
        \begin{eqa}
            S=\sqrt{3}\cdot U\cdot I=\sqrt{3}\cdot U \cdot \sqrt{3}\cdot \frac{U}{Z}=3\cdot \frac{U^2}{Z}
        \end{eqa}
        
        Accordingly, the power turnover at the same impedances in a delta connection is three times greater than in a star connection.  
        }
    
    \b{
        \begin{itemize}
            \item Power is calculated using the equation $P=U\cdot I$
            \item In three-phase current, the power must be calculated for each phase
            \item In a symmetrical case, the currents and voltage in the phases are equal, so the following applies:
        \end{itemize}
        \begin{eqa}
            S&=3\cdot I_{\mathrm{str}}\cdot U_{\Stern}=3\cdot I_{\mathrm{str}}\cdot \frac{U_{\Dreieck}}{\sqrt{3}} \notag \\
            &=3\cdot I_{\Dreieck}\cdot U_{\Dreieck}=3\cdot \frac{I_{\mathrm{str}}}{\sqrt{3}}\cdot U_{\Dreieck} \notag \\
            &=\sqrt{3}\cdot U_{\Dreieck}\cdot I_{\mathrm{str}}
        \end{eqa}
        \begin{itemize}
            \item Fundamentally, it makes no difference which circuit is used for the calculation; only the appropriate values need to be used.
        \end{itemize}
    }
\end{frame}

\begin{frame}
    \b{
        \ftx{Power in a symmetrical three-phase AC network}
        
        \begin{itemize}
            \item With the same impedance, the power in the delta connection is three times greater than in the star connection:
        \end{itemize}
        
        Star connection:
        \begin{eqa}
            S=\sqrt{3}\cdot U\cdot I=\sqrt{3}\cdot U \cdot \frac{U_{\Stern}}{Z}=\frac{U^2}{Z}
        \end{eqa}
        Triangular connection:
        \begin{eqa}
            S=\sqrt{3}\cdot U\cdot I=\sqrt{3}\cdot U \cdot \sqrt{3}\cdot \frac{U}{Z}=3\cdot \frac{U^2}{Z}
        \end{eqa}
    }
\end{frame}



\begin{frame}
    \ftb{Unbalanced components}
    
    \s{
        The fundamentals of a three-phase system were explained in the previous chapter.
        However, it was always assumed that both the generation and the consumers are symmetrical.
        Now let us consider how the system behaves when asymmetries are present.
        However, for the time being, the generator side remains symmetrical and connected in a star configuration.
        If the generation remains star-connected, there can be three different types of circuits:
        
        \begin{itemize}
            \item Four-wire network in star connection
            \item Three-wire network in star connection
            \item Three-wire network in delta connection
        \end{itemize}
    }
    
    \b{
        \begin{itemize}
            \item Until now, it has been assumed that generators and consumers are symmetrical
            \item Now let us consider how the system behaves when asymmetries are present
            \item However, the generator side is still assumed to be symmetrical and connected in a star configuration
            \item For consumers, a distinction is made between four-wire networks in star configuration, three-wire networks in star configuration, and three-wire networks in delta configuration
        \end{itemize}
    }
\end{frame}

\begin{frame}
    \ftx{Asymmetrical components}
    
    \begin{Key point}{Unbalanced components}
        {\bf Unbalanced components} unlike symmetrical components, they describe the {\bf uneven load on the
            phases} due to different consumers. 
    \end{Key point}
\end{frame}

\begin{frame}
    \ftc{Four-wire network in star connection}
    \s{
        The structure of a four-wire network in star connection is shown as ESB in Fig. \ref{ESBStern}.
        The return conductor is connected to both star points, thus connecting the generator and consumer sides. 
        Because the star voltages are symmetrical (based on the assumption that the generator side is symmetrical), the consumer voltages are also voltage-symmetrical.
        However, the currents differ from each other due to the different impedances.
        As a result, the current in the return conductor is usually not equal to zero.
        
        \begin{eqa}
            \underline{I}_\mathrm{N}=\underline{I}_{\mathrm{L1}}+\underline{I}_{\mathrm{L2}}+\underline{I}_{\mathrm{L3}}
        \end{eqa}
        
        Logically, the power can then no longer be calculated as shown in formula \ref{LeistungDrehstrom}.
        Each conductor must now be calculated individually and finally added together.

        \begin{eqa}
            \underline{S}=\underline{I}^*_{\mathrm{L1}}\cdot \underline{U}_{\mathrm{L1}}+\underline{I}^*_{\mathrm{L2}}\cdot \underline{U}_{\mathrm{L2}}+\underline{I}^*_{\mathrm{L3}}\cdot \underline{U}_{\mathrm{L3}}
        \end{eqa}
    }
    \b{
        \begin{itemize}
            \item The structure is the same as for the ESB from the chapter on star-connected systems
            \item The return conductor connects the star points on the generator and consumer sides
            \item If the generator side remains symmetrical, the star voltages remain symmetrical
            \item The currents adjust according to the different impedances
            \item The following applies to the current in the return conductor:
        \end{itemize}
        \begin{eqa}
            \underline{I}_\mathrm{N}=\underline{I}_{\mathrm{L1}}+\underline{I}_{\mathrm{L2}}+\underline{I}_{\mathrm{L3}}
        \end{eqa}
        \begin{itemize}
            \item The power is calculated using the familiar formula
            \item In the unbalanced case, it must be calculated individually for each conductor.
        \end{itemize}
        \begin{eqa}
            \underline{S}=\underline{I}^*_{\mathrm{L1}}\cdot \underline{U}_{\mathrm{L1}}+\underline{I}^*_{\mathrm{L2}}\cdot \underline{U}_{\mathrm{L2}}+\underline{I}^*_{\mathrm{L3}}\cdot \underline{U}_{\mathrm{L3}}
        \end{eqa}
    }
\end{frame}

\begin{frame}
    \renewcommand{\figurename}{Figure}
    \ftc{Three-wire network in star connection}
    
    \s{
        In principle, the same conditions apply in a three-wire star connection as in a four-wire connection.
        The difference is that the current cannot flow back through the return conductor and therefore has an effect on the other conductors (see Fig. \ref{ESBDreileiterDreieck}).
        The same applies to the voltage.         
        Since there is no voltage drop in the return conductor, this affects the voltage across the loads.
        A voltage difference arises between the two star points, which must first be calculated using the voltages across the loads.
        Two methods can be used to determine the required values.         
        For the first method, the impedances are converted mathematically into a delta connection so that they can be calculated easily using known methods.
        This is because in a delta connection, the voltages across the impedances are equal to the conductor voltages.
        The outer conductor currents can also be determined relatively easily.
    }
    
    \b{
        Simplified ESB: Three-wire network in star connection with unbalanced consumers
    }
    
    \fu{
        \input{Tikz/src/Drehstrom/Stern_Dreieck_vereinfachtes_ESB}
    }{{\bf Simplified ESB.} Three-wire network in star connection with unbalanced loads \label{ESBDreileiterDreieck}}
    
\end{frame}

\begin{frame}
    \ftx{Three-wire network in star connection}
    
    
    \s{
    
    
    
        With the alternative solution, equations can be set up using Ohm's and Kirchhoff's laws with the known values of the system.
        The following equation can be set up for the differential voltage between the two star points:
        
        \begin{eqa}
            \underline{U}_{VQ}=\frac{\frac{\underline{U}_{\mathrm{L1}}}{\underline{Z}_1}+\frac{\underline{U}_{\mathrm{L2}}}{\underline{Z}_2}+\frac{\underline{U}_{\mathrm{L3}}}{\underline{Z}_3}}{\frac{1}{\underline{Z}_1}+\frac{1}{\underline{Z}_2}+\frac{1}{\underline{Z}_3}}=\frac{\underline{U}_{\mathrm{L1}}\cdot \underline{Z}_2\cdot \underline{Z}_3+\underline{U}_{\mathrm{L2}}\cdot \underline{Z}_3\cdot \underline{Z}_1+\underline{U}_{\mathrm{L3}}\cdot \underline{Z}_1\cdot \underline{Z}_2}{\underline{Z}_1\cdot \underline{Z}_2+\underline{Z}_2\cdot \underline{Z}_3+\underline{Z}_3\cdot \underline{Z}_1}
        \end{eqa}
        
        As is well known, the power is calculated from the sum of the powers of the individual phases.
        However, it is important to note that the calculation is not based on the outer conductor voltage, but on the voltages across the impedances:
        
        \begin{eqa}
            \underline{S}=\underline{I}^*_{\mathrm{L1}}\cdot \underline{U}_{\mathrm{Z1}}+\underline{I}^*_{\mathrm{L2}}\cdot \underline{U}_{\mathrm{Z2}}+\underline{I}^*_{\mathrm{L3}}\cdot \underline{U}_{\mathrm{Z3}}
        \end{eqa}
        
        
        Theoretically, this formula could be used to perform the calculation.
        However, the measurements do not provide the values required for this formula.
        (How measurements are performed in a three-phase system will be explained in a later chapter.)
        If the formula for power is rearranged as follows, it can be seen that the sum of the conductor currents must be zero, and thus the last term is omitted.        
        
        \begin{eqa}
            \underline{S}&=\underline{I}^*_{\mathrm{L1}}\cdot \underline{U}_{\mathrm{L1}}+\underline{I}^*_{\mathrm{L2}}\cdot \underline{U}_{\mathrm{L2}}+\underline{I}^*_{\mathrm{L3}}\cdot \underline{U}_{\mathrm{L3}}-\underline{U}_{\mathrm{VQ}}\cdot (\underline{I}^*_{\mathrm{L1}}+\underline{I}^*_{\mathrm{L2}}+\underline{I}^*_{\mathrm{L3}}) \notag \\
            &=\underline{I}^*_{\mathrm{L1}}\cdot \underline{U}_{\mathrm{L1}}+\underline{I}^*_{\mathrm{L2}}\cdot \underline{U}_{\mathrm{L2}}+\underline{I}^*_{\mathrm{L3}}\cdot \underline{U}_{\mathrm{L3}}
        \end{eqa}
        
        Here you can see that this formula for power is the same as in a four-wire network in star connection.
        However, a simplification can be made here, as the node rule can be used to express one current in terms of the other two currents:
         
        \begin{eqa}
            \underline{S}&=\underline{I}^*_{\mathrm{\mathrm{L1}}}\cdot \underline{U}_{\mathrm{L1}}+\underline{I}^*_{\mathrm{L2}}\cdot \underline{U}_{\mathrm{L2}}+(-\underline{I}^*_{\mathrm{L1}}-\underline{I}^*_{\mathrm{L2}})\cdot \underline{U}_{\mathrm{L3}} \notag \\
            &=(\underline{U}_{\mathrm{L1}}-\underline{U}_{\mathrm{L2}})\cdot \underline{I}^*_{\mathrm{L1}}+(\underline{U}_{\mathrm{L2}}-\underline{U}_{\mathrm{L3}})\cdot \underline{I}^*_{\mathrm{L2}} \notag \\
            &=\underline{U}_{\mathrm{L3L1}}\cdot \underline{I}^*_{\mathrm{L1}}+ \underline{U}_{\mathrm{L2L3}}\cdot \underline{I}^*_{\mathrm{L2}}
        \end{eqa}
        
    }
    
    \b{
        \begin{itemize}
            \item Similar conditions as for the four-wire network in star connection, but without a return conductor
            \item This prevents the current from flowing back via the return conductor and therefore affects the other phases
            \item Accordingly, the voltage drop between the star points and across the loads must also be calculated
            \item The values can be determined using two methods:
        \end{itemize}
        \begin{itemize}
            \item [1.] Convert impedances mathematically into a delta connection and calculate them using the known methods
            \item Because: In a delta connection, the voltages across the impedances are equal to the conductor voltages
        \end{itemize}
        
    }
    
\end{frame}

\begin{frame}
    \ftx{Three-wire network in star connection}
    
    \b{
        \begin{itemize}
            \item [2.] Set up equations using Ohm's and Kirchhoff's laws
            \item The following equation can be set up for the difference voltage between the two star points:        
        \end{itemize}
        \begin{eqa}
            \underline{U}_{VQ}=\frac{\frac{\underline{U}_{\mathrm{L1}}}{\underline{Z}_1}+\frac{\underline{U}_{\mathrm{L2}}}{\underline{Z}_2}+\frac{\underline{U}_{\mathrm{L3}}}{\underline{Z}_3}}{\frac{1}{\underline{Z}_1}+\frac{1}{\underline{Z}_2}+\frac{1}{\underline{Z}_3}}=\frac{\underline{U}_{\mathrm{L1}}\cdot \underline{Z}_2\cdot \underline{Z}_3+\underline{U}_{\mathrm{L2}}\cdot \underline{Z}_3\cdot \underline{Z}_1+\underline{U}_{\mathrm{L3}}\cdot \underline{Z}_1\cdot \underline{Z}_2}{\underline{Z}_1\cdot \underline{Z}_2+\underline{Z}_2\cdot \underline{Z}_3+\underline{Z}_3\cdot \underline{Z}_1}
        \end{eqa}
        \begin{itemize}
            \item The power is calculated in each phase; it is important to measure the voltage across the consumers.
        \end{itemize}
        \begin{eqa}
            \underline{S}=\underline{I}^*_{\mathrm{L1}}\cdot \underline{U}_{\mathrm{Z1}}+\underline{I}^*_{\mathrm{L2}}\cdot \underline{U}_{\mathrm{Z2}}+\underline{I}^*_{\mathrm{L3}}\cdot \underline{U}_{\mathrm{Z3}}
        \end{eqa}
    }
\end{frame}

\begin{frame}
    \ftx{Three-wire network in star connection}
    
    \b{
        \begin{itemize}
            \item The equation for calculating power can be rearranged and simplified
        \end{itemize}
        \begin{eqa}
            \underline{S}&=\underline{I}^*_{\mathrm{L1}}\cdot \underline{U}_{\mathrm{L1}}+\underline{I}^*_{\mathrm{L2}}\cdot \underline{U}_{\mathrm{L2}}+\underline{I}^*_{\mathrm{L3}}\cdot \underline{U}_{\mathrm{L3}}-\underline{U}_{\mathrm{VQ}}\cdot (\underline{I}^*_{\mathrm{L1}}+\underline{I}^*_{\mathrm{L2}}+\underline{I}^*_{\mathrm{L3}}) \notag \\
            &=\underline{I}^*_{\mathrm{L1}}\cdot \underline{U}_{\mathrm{L1}}+\underline{I}^*_{\mathrm{L2}}\cdot \underline{U}_{\mathrm{L2}}+\underline{I}^*_{\mathrm{L3}}\cdot \underline{U}_{\mathrm{L3}}
        \end{eqa}
        \begin{itemize}
            \item Using the node rule, one of the currents can in turn be expressed by the other two
        \end{itemize}
        \begin{eqa}
            \underline{S}&=\underline{I}^*_{\mathrm{\mathrm{L1}}}\cdot \underline{U}_{\mathrm{L1}}+\underline{I}^*_{\mathrm{L2}}\cdot \underline{U}_{\mathrm{L2}}+(-\underline{I}^*_{\mathrm{L1}}-\underline{I}^*_{\mathrm{L2}})\cdot \underline{U}_{\mathrm{L3}} \notag \\
            &=(\underline{U}_{\mathrm{L1}}-\underline{U}_{\mathrm{L2}})\cdot \underline{I}^*_{\mathrm{L1}}+(\underline{U}_{\mathrm{L2}}-\underline{U}_{\mathrm{L3}})\cdot \underline{I}^*_{\mathrm{L2}} \notag \\
            &=\underline{U}_{\mathrm{L3L1}}\cdot \underline{I}^*_{\mathrm{L1}}+ \underline{U}_{\mathrm{L2L3}}\cdot \underline{I}^*_{\mathrm{L2}}
        \end{eqa}
    }
\end{frame}

\begin{frame}
    \ftc{Three-wire network in delta connection}
    
    \s{
    
        In a delta connection, there is no neutral point, which is why the delta connection can only be implemented in a three-wire network.
        In the delta connection, the outer conductor voltages are also symmetrical due to the symmetrical feed.
        The currents adjust according to the impedances and can be calculated as already known:
        
        \begin{eqa}
            \underline{I}_{\mathrm{L1}}=\underline{I}_{\mathrm{L1L2}}-\underline{I}_{\mathrm{L3L1}} \\
            \underline{I}_{\mathrm{L2}}=\underline{I}_{\mathrm{L2L3}}-\underline{I}_{\mathrm{L1L2}} \\
            \underline{I}_{\mathrm{L3}}=\underline{I}_{\mathrm{L3L1}}-\underline{I}_{\mathrm{L2L3}}
        \end{eqa}
        
        
        The sum of the outer conductor currents must, of course, always be zero, even with unbalanced currents.
        This results in the same calculation for the power calculation as for the star connection in a three-wire network and therefore also the same transformations.        
        \begin{eqa}
            \underline{S}&=\underline{I}^*_{\mathrm{L1L2}}\cdot \underline{U}_{\mathrm{L1L2}}+\underline{I}^*_{\mathrm{L2L3}}\cdot \underline{U}_{\mathrm{L2L3}}+\underline{I}^*_{\mathrm{L3L1}}\cdot \underline{U}_{\mathrm{L3L1}} \notag \\ 
            &=\underline{I}^*_{\mathrm{L1}}\cdot \underline{U}_{\mathrm{L1}}+\underline{I}^*_{\mathrm{L2}}\cdot \underline{U}_{\mathrm{L2}}+\underline{I}^*_{\mathrm{L3}}\cdot \underline{U}_{\mathrm{L3}} \\ \notag \\
            \underline{S}&=\underline{I}^*_{\mathrm{L1}}\cdot \underline{U}_{\mathrm{L1}}+\underline{I}^*_{\mathrm{L2}}\cdot \underline{U}_{\mathrm{L2}}+(-\underline{I}^*_{\mathrm{L1}}-\underline{I}^*_{\mathrm{L2}})\cdot \underline{U}_{L3} \notag \\
            &=(\underline{U}_{\mathrm{L1}}-\underline{U}_{\mathrm{L2}})\cdot \underline{I}^*_{\mathrm{L1}}+(\underline{U}_{\mathrm{L2}}-\underline{U}_{\mathrm{L3}})\cdot \underline{I}^*_{\mathrm{L2}} \notag \\
            &=\underline{U}_{\mathrm{L3L1}}\cdot \underline{I}^*_{\mathrm{L1}}+ \underline{U}_{\mathrm{L2L3}}\cdot \underline{I}^*_{\mathrm{L2}}
        \end{eqa}
    }
    
    \b{
        \begin{itemize}
            \item There is no neutral point available in the delta connection, so it can only be implemented in a three-wire system.
            \item Due to the symmetrical feed, the outer conductor voltages are symmetrical.
            \item The currents adjust to the respective impedance:
        \end{itemize}
        \begin{eqa}
            \underline{I}_{\mathrm{L1}}=\underline{I}_{\mathrm{L1L2}}-\underline{I}_{\mathrm{L3L1}} \\
            \underline{I}_{\mathrm{L2}}=\underline{I}_{\mathrm{L2L3}}-\underline{I}_{\mathrm{L1L2}} \\
            \underline{I}_{\mathrm{L3}}=\underline{I}_{\mathrm{L3L1}}-\underline{I}_{\mathrm{L2L3}}
        \end{eqa}
    }
    
\end{frame}

\begin{frame}
    \ftx{Three-wire network in delta connection}
    \b{
    
        \begin{itemize}
            \item The sum of the external conductor currents must always be zero, even with unbalanced currents.
            \item This results in the same power calculation as for the star connection in a three-wire network and therefore also the same transformations.
        \end{itemize}
        \begin{eqa}
            \underline{S}&=\underline{I}^*_{\mathrm{L1L2}}\cdot \underline{U}_{\mathrm{L1L2}}+\underline{I}^*_{\mathrm{L2L3}}\cdot \underline{U}_{\mathrm{L2L3}}+\underline{I}^*_{\mathrm{L3L1}}\cdot \underline{U}_{\mathrm{L3L1}} \notag \\ 
            &=\underline{I}^*_{\mathrm{L1}}\cdot \underline{U}_{\mathrm{L1}}+\underline{I}^*_{\mathrm{L2}}\cdot \underline{U}_{\mathrm{L2}}+\underline{I}^*_{\mathrm{L3}}\cdot \underline{U}_{\mathrm{L3}} \\ \notag \\
            \underline{S}&=\underline{I}^*_{\mathrm{L1}}\cdot \underline{U}_{\mathrm{L1}}+\underline{I}^*_{\mathrm{L2}}\cdot \underline{U}_{\mathrm{L2}}+(-\underline{I}^*_{\mathrm{L1}}-\underline{I}^*_{\mathrm{L2}})\cdot \underline{U}_{L3} \notag \\
            &=(\underline{U}_{\mathrm{L1}}-\underline{U}_{\mathrm{L2}})\cdot \underline{I}^*_{\mathrm{L1}}+(\underline{U}_{\mathrm{L2}}-\underline{U}_{\mathrm{L3}})\cdot \underline{I}^*_{\mathrm{L2}} \notag \\
            &=\underline{U}_{\mathrm{L3L1}}\cdot \underline{I}^*_{\mathrm{L1}}+ \underline{U}_{\mathrm{L2L3}}\cdot \underline{I}^*_{\mathrm{L2}}
        \end{eqa}
        
    }
\end{frame}

\newpage